\documentclass[11pt]{article}
\usepackage[margin=1in]{geometry}
\usepackage{amsmath}
\usepackage{amssymb}
\usepackage{enumitem}
\usepackage[hidelinks]{hyperref}

\title{CSCE 355: JFLAP5 --- Pushdown Automata}
\author{}
\date{}

\renewcommand{\labelitemi}{$\bullet$}
\newcommand{\fn}[1]{\texttt{\detokenize{#1}}}
\newcommand{\tr}[3]{\texttt{#1, #2 ; #3}}

\begin{document}
\maketitle

\section*{Overview}

In this assignment you will build five pushdown automata (PDAs) in JFLAP. Each one is checked
automatically. The assignment is worth \textbf{50 points}. Throughout, $\mathbb{N} = \{0, 1, 2,
\ldots\}$ includes 0, and $\lambda$ is the empty string.

\subsection*{Getting started}

\begin{itemize}
  \item Read the JFLAP PDA tutorial:
        \url{https://www.jflap.org/tutorial/pda/construct/index.html}.
  \item JFLAP is a free \texttt{.jar} file from \url{https://www.jflap.org}. You need Java 8 or
        higher. Start it with \texttt{java -jar JFLAP7\_1.jar} (use the exact name of the file you
        downloaded; on macOS and Linux, upper and lower case matter).
  \item Unzip this assignment folder. (It is a good idea to append your username to the folder
        name, e.g.\ \fn{JFLAP5_jsmith}.)
  \item In JFLAP use \textbf{File $\to$ Open} to open one of the provided \texttt{.jff} files.
        Each one is an empty \emph{Pushdown Automaton}. Build your PDA, then \textbf{File $\to$
        Save} it \emph{under the same file name}. If you ever start a new PDA instead, choose
        \textbf{Pushdown Automaton} and then \textbf{Multiple Character Input}.
\end{itemize}

\begin{center}
\begin{tabular}{|c|l|l|c|}
\hline
\textbf{\#} & \textbf{Your PDA (edit this file)} & \textbf{Test strings (do not edit)} & \textbf{Points}\\
\hline
1 & \fn{problem1_axbycx.jff}         & \fn{problem1_axbycx_tests.txt}         & 10 \\
2 & \fn{problem2_not_ww.jff}         & \fn{problem2_not_ww_tests.txt}         & 20 \\
3 & \fn{problem3_0n1n_or_0n12n.jff}  & \fn{problem3_0n1n_or_0n12n_tests.txt}  & 10 \\
4 & \fn{problem4_aibjck.jff}         & \fn{problem4_aibjck_tests.txt}         & 5 \\
5 & \fn{problem5_anbm.jff}           & \fn{problem5_anbm_tests.txt}           & 5 \\
\hline
\end{tabular}
\end{center}

\section*{How these PDAs work}

\begin{itemize}
  \item \textbf{Transitions.} JFLAP writes a transition as \tr{read}{pop}{push}. It reads
        \texttt{read} from the input, pops the string \texttt{pop} from the top of the stack, and
        pushes the string \texttt{push}, whose first character ends up on top. Any of the three
        may be $\lambda$: leave the field blank and JFLAP shows $\lambda$. For example,
        \tr{a}{Z}{aZ} reads an \texttt{a} when \texttt{Z} is on top and leaves \texttt{a} on top
        of \texttt{Z}.
  \item \textbf{The stack starts with \texttt{Z}}, the bottom-of-stack marker, and nothing else.
        You may use any other characters as stack symbols.
  \item \textbf{Accept by final state.} A string is accepted when the PDA can read all of it and
        be in a final state; what is left on the stack does not matter. When JFLAP asks
        ``Accept by'', choose \textbf{Final State}.
  \item \textbf{Nondeterminism is allowed} and often needed: a state may have several
        transitions that apply at the same time, and $\lambda$ transitions.
\end{itemize}

\subsection*{Rules for every PDA}
These are checked automatically, before any test string is run. A PDA that breaks one of them
receives \textbf{zero credit} for that problem.
\begin{itemize}
  \item Exactly one initial state, and at least one final state.
  \item The input part of every label uses only symbols of the problem's alphabet. It may read
        several symbols at once (that is what ``Multiple Character Input'' allows) or none
        ($\lambda$). Do not type a comma-separated list such as \texttt{a,b}: JFLAP reads that as
        the three-character string \texttt{a,b}. Draw a separate transition for each choice.
\end{itemize}

\section*{Problems}
In every test file, \textbf{the first 5 strings should reject and the last 5 should accept.}

\subsection*{Problem 1 (10 points)}
Construct a PDA recognizing
\[
  L = \{\, a^x b^y c^x \;:\; x, y \in \mathbb{N} \,\}
\]
over the alphabet $\{a, b, c\}$: some \texttt{a}s, then any number of \texttt{b}s, then exactly
as many \texttt{c}s as \texttt{a}s. For example, \texttt{abc}, \texttt{bbb} and
\texttt{aaacc\-c} are accepted; \texttt{cba} and \texttt{ccc} are not. The empty string should
\textbf{accept} ($x = y = 0$).

To get you started, the start state can read the \texttt{a}s with these two self-loops:
\tr{a}{Z}{aZ} and \tr{a}{a}{aa}.

\subsection*{Problem 2 (20 points)}
Design a PDA for the set of all strings of \texttt{a}s and \texttt{b}s that are \emph{not} of the
form $ww$, that is, not equal to any string written twice:
\[
  L = \{\, s \in \{a,b\}^* \;:\; s \ne ww \text{ for every } w \in \{a,b\}^* \,\}
\]
For example, \texttt{a}, \texttt{abbaaaba} and \texttt{bbbbabbb} are accepted; \texttt{abab}
($w = $ \texttt{ab}) and \texttt{bbbbbbbb} are not. The empty string should \textbf{reject}
(it is $ww$ with $w = \lambda$).

\emph{Hint 1:} write down some strings that are ``almost'' of the form $ww$ and see what they
have in common. \emph{Hint 2:} look at the shape
\texttt{xxxxxx\,a\,xxxxxx\,yy\,b\,yy}, where each \texttt{x} and \texttt{y} stands for any
symbol.

\subsection*{Problem 3 (10 points)}
Construct a PDA recognizing
\[
  L = \{\, 0^n 1^n \;:\; n \in \mathbb{N},\ n \ge 1 \,\} \;\cup\; \{\, 0^n 1^{2n} \;:\; n \in
  \mathbb{N},\ n \ge 1 \,\}
\]
over the alphabet $\{0, 1\}$. For example, \texttt{0011}, \texttt{011} and \texttt{001111} are
accepted; \texttt{00111} and \texttt{010101} are not. The empty string should \textbf{reject}
($n \ge 1$).

\subsection*{Problem 4 (5 points)}
Construct a PDA recognizing
\[
  L = \{\, a^i b^j c^k \;:\; i, j, k \in \mathbb{N},\ j = i + k \,\}
\]
over the alphabet $\{a, b, c\}$: some \texttt{a}s, then some \texttt{b}s, then some
\texttt{c}s, with as many \texttt{b}s as \texttt{a}s and \texttt{c}s together. For example,
\texttt{abbc}, \texttt{aabb} and \texttt{bbcc} are accepted; \texttt{abc} and \texttt{acb} are
not. The empty string should \textbf{accept}.

\subsection*{Problem 5 (5 points)}
Construct a PDA recognizing
\[
  L = \{\, a^n b^m \;:\; n, m \in \mathbb{N},\ n \le m \le 2n \,\}
\]
over the alphabet $\{a, b\}$: some \texttt{a}s followed by at least as many \texttt{b}s, but no
more than twice as many. For example, \texttt{ab}, \texttt{abb} and \texttt{aabbb} are accepted;
\texttt{aab} and \texttt{abbb} are not. The empty string should \textbf{accept}.

\section*{Testing your PDAs}

\subsection*{In JFLAP}
\begin{enumerate}
  \item Open your PDA (it needs an initial state to be tested).
  \item Choose \textbf{Input $\to$ Multiple Run}. When asked ``Accept by'', choose
        \textbf{Final State}.
  \item Click \textbf{Load Inputs} and select the problem's test file.
  \item Click \textbf{Run Inputs}. The first 5 rows should say Reject and the last 5 Accept.
        For a very nondeterministic PDA (Problem 2 especially) JFLAP may ask whether to continue
        after 10000 configurations; answer Yes.
\end{enumerate}

\subsubsection*{Testing the empty string}
The test files do not contain the empty string. In Multiple Run, put the cursor on the last row
and press Enter to add a blank line (an empty input), then run again. The empty string should
accept in Problems 1, 4 and 5, and reject in Problems 2 and 3.

\subsection*{From the command line}
This folder contains \texttt{test\_pda.py}, the same script that is used to grade you (Python 3
is required). From inside this folder run
\begin{center}
\texttt{python3 test\_pda.py}
\end{center}
to check all five problems, or \texttt{python3 test\_pda.py 2} to check only Problem 2. The
report tells you, for each problem, whether the PDA follows the rules above, which test strings
are wrong, a shortest string your PDA gets wrong (if any), and your score. Run it as often as you
like.

\section*{Grading}

Each problem earns its points only if \emph{all} of the following hold:
\begin{itemize}
  \item The PDA follows the rules above.
  \item Every string in the problem's test file gets the expected result, and so does the empty
        string.
  \item The PDA gets the right answer on \emph{every} string up to length 10 over $\{a,b\}$,
        length 7 over $\{a,b,c\}$, and length 12 over $\{0,1\}$, and on some longer strings.
        The provided strings are only examples, so a PDA that passes them but is wrong elsewhere
        earns no credit. When this happens the report shows a shortest string on which your PDA
        is wrong.
\end{itemize}

\textbf{Do not edit the test files or \texttt{test\_pda.py}.} They are replaced with the
original versions before grading, so changing them can only make your own testing misleading.

\section*{Submission}

\begin{enumerate}
  \item Open \textbf{YOURUSERNAME.txt} and put your university username as the \emph{first line}
        of the file (only your username, nothing else on that line). This is how your submission
        is identified.
  \item Make sure each of your five PDAs is saved in its provided \texttt{.jff} file. Do not
        rename any of the provided files.
  \item Run \texttt{python3 test\_pda.py} one last time to make sure everything is as you expect.
  \item Zip this entire folder and submit the zip as one file.
\end{enumerate}

\end{document}
